\introduction

Актуальность темы исследования обусловлена тем, что молочное животноводство является одной из ключевых отраслей агропромышленного комплекса: по данным ФАО, мировое производство коровьего молока превышает 750~млн тонн в год и продолжает расти со среднегодовым темпом 1,5--2~процента~\cite{PrecisionLivestock}. Эффективность управления молочной фермой напрямую влияет на рентабельность производства: своевременное выявление мастита снижает экономические потери на 30--50~процентов, а оптимизация кормления позволяет повысить надои на 5--15~процентов~\cite{DairyPrecision}.

Концепция точного животноводства (precision livestock farming), предложенная Д.~Беркмансом в 2008~году, предполагает непрерывный автоматизированный мониторинг состояния животных с использованием датчиков и информационных технологий~\cite{PrecisionLivestock}. Современные молочные фермы оснащаются роботизированными доильными установками, датчиками активности и жвачки, автоматическими кормушками, генерирующими значительные объёмы данных --- до 250\,000 записей в сутки для стада в 500~голов~\cite{DairyPrecision}. Без специализированной информационной системы оперативная обработка этих данных невозможна, что приводит к запоздалому выявлению заболеваний, неоптимальному кормлению и потерям продукции.

Существующие системы управления молочными фермами (Lely Horizon, DeLaval DelPro, SCR Heatime, 1С: Управление фермой) обладают существенными ограничениями: монолитная архитектура, затрудняющая расширение; отсутствие встроенных средств предиктивной аналитики; привязка к оборудованию одного производителя; закрытый API, не позволяющий внедрять собственные модели машинного обучения. Поэтому актуальна задача создания информационной системы, объединяющей полный цикл управления стадом, интеграцию с роботизированным оборудованием и модули машинного обучения в единой открытой платформе.

Объектом исследования является процесс автоматизации управления молочной фермой. Предметом исследования --- методы и средства разработки информационных систем управления молочными фермами с интеграцией роботизированного доильного оборудования и предиктивной аналитикой на основе моделей машинного обучения. Цель работы --- разработка информационной системы управления молочной фермой с интеграцией роботизированного доильного оборудования Lely и предиктивной аналитикой на основе моделей машинного обучения.

Для достижения поставленной цели необходимо решить следующие задачи:
\begin{enumerate}
    \item провести анализ предметной области и существующих решений, на основании которого сформулировать техническое задание и обосновать выбор стека технологий;
    \item спроектировать трёхзвенную архитектуру и модель данных базы данных, содержащую более 35 таблиц с расширением TimescaleDB для работы с временными рядами;
    \item реализовать серверный API на Rust/Axum с аутентификацией, авторизацией, middleware безопасности и интеграцией с облачным API Lely Horizon;
    \item разработать клиентское веб-приложение и модуль машинного обучения с 11 моделями для предиктивной аналитики;
    \item реализовать серверные алгоритмы прогнозирования, ансамблевый метод и объяснимость прогнозов на основе SHAP;
    \item провести функциональное и нагрузочное тестирование системы, а также разработать конвейер CI/CD и мониторинг.
\end{enumerate}

В работе использованы следующие методы исследования: системный анализ предметной области, проектирование реляционных баз данных с использованием расширенной модели ``сущность--связь'' и нормализации, методы машинного обучения (градиентный бустинг XGBoost с квантильной регрессией, кластеризация KMeans, обнаружение аномалий Isolation Forest), анализ временных рядов (экспоненциальное сглаживание Хольта--Винтерса, ряды Фурье, ARIMA, линейная регрессия), а также методы программной инженерии (контейнеризация Docker, непрерывная интеграция и доставка, мониторинг на базе стека Grafana LGTM). Качество моделей прогнозирования оценивается по метрикам MAPE и RMSE на тестовых выборках с хронологическим разбиением. Объяснимость прогнозов обеспечивается методом SHAP, позволяющим количественно оценить вклад каждого признака.

Научная новизна работы заключается в следующем:
\begin{enumerate}
    \item предложена архитектура информационной системы, интегрирующая серверные алгоритмы аналитики на Rust с внешним модулем машинного обучения на Python, что обеспечивает работу предиктивной аналитики даже при недоступности ML-сервиса;
    \item разработан ансамблевый метод объединения прогнозов с весами, обратно пропорциональными MAPE, повышающий робастность прогнозирования;
    \item реализована система автоматического сравнения 8 моделей прогнозирования временных рядов с выбором оптимальной для каждого животного индивидуально.
\end{enumerate}

Практическая значимость работы заключается в том, что разработанная информационная система может быть использована на молочных фермах, оснащённых роботизированным доильным оборудованием Lely, для автоматизации процессов учёта и мониторинга. Внедрение системы позволяет повысить эффективность управления стадом за счёт предиктивной аналитики, сократить время на ведение документации и своевременно выявлять заболевания животных на ранних стадиях.

Выпускная квалификационная работа состоит из введения, трёх глав, заключения, списка использованных источников и приложений. В первой главе проведён анализ предметной области: рассмотрена концепция точного животноводства, выполнен обзор существующих решений и сформулировано техническое задание. Во второй главе описаны проектирование и разработка системы: архитектура, модель данных, серверная и клиентская части, модуль машинного обучения, серверная аналитика, мониторинг и развёртывание. В третьей главе представлены результаты тестирования и демонстрация работы системы. В заключении сформулированы основные результаты работы и перспективы дальнейшего развития.
